\subsection{Examples of manifolds}

5.2.1. Let X be a set; there exists on X one and only one manifold structure for which the underlying topological space is discrete; this structure is a pure manifold structure of dimension 0.

5.2.2. Let E be a Banach space. The triple \(c = (\mathrm{E}, \mathrm{Id}_{\mathrm{E}}, \mathrm{E})\) is a chart of E and \(\mathcal{A} = \{c\}\) is an atlas of E, hence defines a pure manifold structure of type E on E; the underlying topology is the given topology on E. In what follows, when we speak of the manifold structure on E, it will always be to the preceding structure that we refer.

In particular, this applies to every finite-dimensional vector space over K, equipped with the unique Hausdorff topology compatible with its vector structure (Esp. Vect. Top., chap. I, § 2, no. 3).

5.2.3. Let X be a manifold and U an open subset of X. There exists on U a manifold structure whose charts are the charts of the manifold X with domain contained in U. This structure is said to be induced by that of X (cf. no. 5.3.4); equipped with this structure, U is called an open submanifold of X.

In particular, every open subset of a Banach space E is equipped with a canonical structure of pure manifold of type E. Let X be a manifold; for a triple (U, \(\varphi\) , E) to be a chart of X, it is necessary and sufficient that U be open in X and that \(\varphi\) be an isomorphism of the open submanifold U of X onto an open submanifold of E.

5.2.4. Let X be a set, and \((X_{i})_{i\in I}\) a covering of X. Suppose that on each \(X_{i}\) there is given a manifold structure such that the following condition is satisfied:

For any \(i\) and \(j\) in I, the set \(\mathbf{X}_{i} \cap \mathbf{X}_{j}\) is open in \(\mathbf{X}_{i}\) and \(\mathbf{X}_{j}\) and the manifold structures induced on \(\mathbf{X}_{i} \cap \mathbf{X}_{j}\) by those of \(\mathbf{X}_{i}\) and \(\mathbf{X}_{j}\) coincide.

There then exists on X one and only one manifold structure such that \(X_{i}\) is an open submanifold of X for every i in I. One says that the manifold X is obtained by gluing together the manifolds \(X_{i}\) .

5.2.5. Let X be a manifold. The set \(X_{n}\) of points x of X such that \(\dim_{x}X=n\ (n\ \text{integer}\geqslant0)\) is an open submanifold of X, which is pure of dimension n.

5.2.6. Let E be a Banach space. One can equip the set \(\mathbf{G}(\mathbf{E})\) of vector subspaces of E admitting a topological supplement with an analytic manifold structure in the following way: for every pair \((\mathrm{F}_0, \mathrm{G}_0) \in \mathbf{G}(\mathbf{E}) \times \mathbf{G}(\mathbf{E})\) such that \(\mathrm{E} = \mathrm{F}_0 \oplus \mathrm{G}_0\) , one denotes by \(\mathrm{U}_{\mathrm{G}_0}\) the set of \(\mathrm{F} \in \mathbf{G}(\mathbf{E})\) admitting \(\mathrm{G}_0\) as supplement, and one defines a bijection \(\varphi_{\mathrm{F}_0, \mathrm{G}_0}\) of \(\mathrm{U}_{\mathrm{G}_0}\) onto the Banach space \(\mathcal{L}(\mathrm{F}_0; \mathrm{G}_0)\) by associating to every \(\mathrm{F} \in \mathrm{U}_{\mathrm{G}_0}\) the map from \(\mathrm{F}_0\) to \(\mathrm{G}_0\) having as graph the subspace F of \(\mathrm{E} = \mathrm{F}_0 \times \mathrm{G}_0\) . The charts \((\mathrm{U}_{\mathrm{G}_0}, \varphi_{\mathrm{F}_0, \mathrm{G}_0}, \mathcal{L}(\mathrm{F}_0; \mathrm{G}_0))\) form an atlas of \(\mathbf{G}(\mathbf{E})\) . Equipped with the manifold structure defined by this atlas, \(\mathbf{G}(\mathbf{E})\) is called the Grassmannian manifold of E.

The topological space G(E) is metrizable. If K is locally compact and E is finite-dimensional, G(E) is compact.

For every integer \(n \geqslant 0\) , the set \(\mathbf{G}_{n}(\mathbf{E})\) (resp. \(\mathbf{G}^{n}(\mathbf{E})\) ) of \(\mathbf{F} \in \mathbf{G}(\mathbf{E})\) of dimension (resp. codimension) \(n\) is open and closed in \(\mathbf{G}(\mathbf{E})\) and is a pure open submanifold of \(\mathbf{G}(\mathbf{E})\) . If \(\mathbf{K} = \mathbf{R}\) or \(\mathbf{C}\) , \(\mathbf{G}_{n}(\mathbf{E})\) is connected for every \(n\) .

The map which associates to each \(F \in \mathbf{G}(\mathbf{E})\) its orthogonal in the dual \(E'\) of E is a morphism of \(\mathbf{G}(\mathbf{E})\) to \(\mathbf{G}(\mathbf{E}')\) which induces an isomorphism of \(\mathbf{G}^{n}(\mathbf{E})\) onto \(\mathbf{G}_{n}(\mathbf{E}')\) for every integer n.

If K = R or C, or if E is finite-dimensional, \(\mathbf{G}_{1}(\mathbf{E})\) is none other than the projective space associated with E, which is thus equipped with an analytic manifold structure.

